\section{Original native fixed-Gauss convex profiles}

Keep an ORIGINAL integral scheme $X$, its ORIGINAL morphism
$f:X\to\operatorname{Spec}k[T]$, and an ORIGINAL point $p\in X$.
Put $S=\mathcal O_{X,p}$, with its actual stalk coefficient map
$\phi_p:k[T]\to S$ induced by $f$, and $t_p=\phi_p(T)$.
The $k$-action is the restriction of THIS map, not a separately supplied
action. Assume this original local ring is Noetherian and has original
parameters $\mathfrak m_S=(\pi,z)$ with $\pi$ prime, $\pi\nmid z$,
$t_p\ne0$, and $\pi\mid t_p$. Fix $0<r<1$.
These are honest local geometric inputs; no global surface existence,
valuation family, field/norm identification, order pair or profile is input.

\begin{lemma}[Canonical original-stalk field points and whole section values]
\label{lem:gauss-native-original-stalk-point-sequence}
\uses{thm:gauss-surface-actual-fixed-radius-points,thm:gauss-surface-actual-fixed-radius-convergence,lem:actual-specified-field-native-valued-generic-point,scheme:valued-point-topology,scheme:relative-section-coefficients,scheme:relative-affine-open-homeomorphism,spectrum:trivial-field-bridge,scheme:relative-valued-points}
\lean{EnvelopeProperty.GaussRegularSurfaceConvexProfiles.stalkCoefficientMap,EnvelopeProperty.GaussRegularSurfaceConvexProfiles.stalkParameter,EnvelopeProperty.GaussRegularSurfaceConvexProfiles.genericFieldMap,EnvelopeProperty.GaussRegularSurfaceConvexProfiles.genericFieldMap_germ,EnvelopeProperty.GaussRegularSurfaceConvexProfiles.nativeWeightedPoint,EnvelopeProperty.GaussRegularSurfaceConvexProfiles.nativeVerticalPoint,EnvelopeProperty.GaussRegularSurfaceConvexProfiles.nativeWeightedPoint_underlying,EnvelopeProperty.GaussRegularSurfaceConvexProfiles.nativeVerticalPoint_underlying,EnvelopeProperty.GaussRegularSurfaceConvexProfiles.nativeWeightedPoint_eval,EnvelopeProperty.GaussRegularSurfaceConvexProfiles.nativeVerticalPoint_eval,EnvelopeProperty.GaussRegularSurfaceConvexProfiles.tendsto_nativeWeightedPoint}
\leanok
\leavevmode\par
Let $K=\operatorname{Frac}S$. Construct the SAME normalized actual points
$w_N,w_{\mathrm{vert}}$ of the frozen real-norm theorem. There are canonical ORIGINAL
native points $q_N,q_{\mathrm{vert}}$ of $X$ over the exact radius-$r$ Gauss base, such
that $q_N\to q_{\mathrm{vert}}$, all underlying scheme points equal the original generic
point, and for EVERY original nonempty open $V$ and original section $a$
its actual section value at $q_N$ or $q_{\mathrm{vert}}$ equals its value at $w_N$ or
$w_{\mathrm{vert}}$ after the canonical ORIGINAL generic germ in $K$.
\end{lemma}
\begin{proof}
The actual original function field of $X$ is the fraction ring of its
actual original stalk $S$. Thus the canonical fraction-ring algebra
isomorphism $K\simeq_S K(X)$ is internally available. Give $K$ the inverse
actual generic-field action and apply the existing specified-field native
generic-point construction. Its whole-section evaluation theorem gives
exactly the asserted equality. The actual stalk/generic-germ square
identifies the image of every original polynomial coefficient with
$\phi_p$ in $S$ followed by its original map into $K$. The frozen theorem's
EVERY-polynomial Gauss equality therefore proves actual membership in the
exact prescribed-base native locus. Its whole-field convergence and the
existing continuity of the specified-field generic-point map prove the
native convergence. Passing to the fixed-base subspace preserves it.
The existing underlying-point theorem gives the original generic point.
Every original nonempty Zariski open contains that point, so all of the
claimed evaluations are defined. No original field or norm square is
assumed: each comes from the actual stalk and generic-germ maps.
\end{proof}

\begin{lemma}[All actual original unit-log expressions have derived profiles]
\label{lem:gauss-native-original-unit-profiles}
\uses{lem:gauss-native-original-stalk-point-sequence,thm:gauss-surface-actual-fixed-radius-convergence,lem:gauss-surface-parameter-positive-vertical-order,def:gauss-surface-same-field-native-orders,lem:gauss-surface-canonical-vertical-leading-field}
\lean{EnvelopeProperty.GaussRegularSurfaceConvexProfiles.nativeScale,EnvelopeProperty.GaussRegularSurfaceConvexProfiles.nativeScale_pos,EnvelopeProperty.GaussRegularSurfaceConvexProfiles.tendsto_nativeScale,EnvelopeProperty.GaussRegularSurfaceConvexProfiles.exists_unit_profiles}
\leanok
\leavevmode\par
For any ORIGINAL nonempty open $V$ and finite original units
\[
 a_i\in\Gamma(X,V)^\times,
\]
internally obtain their true vertical orders
$A_i$ and true canonical leading-residue orders $B_i$. With the SAME actual
parameter orders $e>0,c$ and
\[
 \varepsilon_N=-\log r/d_N(t_p)>0,
\]
eventually
\[
 \log|a_i(q_N)|=\log|a_i(q_{\mathrm{vert}})|
       -\varepsilon_N\bigl(B_i-(A_i/e)c\bigr)
\]
simultaneously for EVERY index $i$. Thus EVERY rational affine unit-log
expression $L_j$ has actual eventual profile
\[
 L_j(q_N)=b_j+\varepsilon_N s_j,\quad
 b_j=L_j(q_{\mathrm{vert}}),\quad
 s_j=-\sum_i j_i\bigl(B_i-(A_i/e)c\bigr).
\]
The orders, profiles, coefficient triangles and nonzero evaluations are
conclusions, not inputs.
\end{lemma}
\begin{proof}
The original generic germ of each original unit is nonzero because the
original integral scheme injects every original nonempty open section ring
into its actual function field. Transport through the canonical isomorphism
of the preceding lemma and apply the frozen all-fraction log-profile
producer to these actual nonzero elements of $K$. Its $e,c$ are selected
once from the SAME original $t_p$, independently of $i$; its $A_i,B_i$ are
native vertical and canonical leading-residue orders. Finiteness of the
index set combines the actual eventual tails. The preceding whole-section
value equality identifies the actual original section logarithms without a
norm-matching premise. Taking the genuine rational linear combination and
adding its actual constant proves the asserted affine profile. Values at
$q_{\mathrm{vert}}$ give exactly $b_j$, including every nonzero vertical intercept.
\end{proof}

\begin{theorem}[SOURCE convexity produces the actual active-index quotient]
\label{thm:gauss-native-convex-active-profile}
\uses{lem:gauss-native-original-unit-profiles,scheme:native-regular-unit-local-psh,analysis:actual-finite-affine-active-max,lem:gauss-native-original-stalk-point-sequence}
\lean{EnvelopeProperty.GaussRegularSurfaceConvexProfiles.exists_convex_active_profile}
\leanok
\leavevmode\par
Let $U$ be an ORIGINAL native domain over the radius-$r$ Gauss base, let
$u$ be an ORIGINAL extended-real function on that same native space, and
assume the literal source predicate
$\operatorname{RelativeSchemeValuedPoint.IsConvexPLOn}(U,u)$ and $q_{\mathrm{vert}}\in U$.
Internally obtain its original local finite unit-log maximum
$u=\max_{j\in I}L_j$ near $q_{\mathrm{vert}}$. Put
\[
 I_0=\{j\in I:b_j=\max_{h\in I}b_h\},\qquad
 D=\max_{j\in I_0}s_j.
\]
Then $I_0$ is genuinely nonempty, and for the SAME actual original native
sequence one has eventually
\[
 u(q_N)\in\mathbf R,\qquad
 u(q_N)=u(q_{\mathrm{vert}})+\varepsilon_ND,\qquad
 \frac{u(q_N)-u(q_{\mathrm{vert}})}{\varepsilon_N}=D.
\]
The true quotient tends to $D$. No supplied list of units, active indices,
profile, slope, derivative, norm equality, harmonic predicate or positivity
conclusion occurs among the inputs.
\end{theorem}
\begin{proof}
Apply the literal source definition at the actual $q_{\mathrm{vert}}$. It supplies an
actual open analytic neighborhood $W$ containing $q_{\mathrm{vert}}$, an original Zariski
open $V$, actual units $a_i$, a nonempty finite rational index set $I$ and
the original maximum identity on $W$. Since $q_N\to q_{\mathrm{vert}}$, the SAME actual
$q_N$ eventually belongs to $W$. Its underlying generic point also belongs
to $V$. Apply the preceding unit-profile lemma to those internally supplied
units and combine its finite tails with this genuine neighborhood tail.

The actual finite-real maximum has a nonempty active set $I_0$. The existing
finite-affine-maximum theorem supplies $\delta>0$ such that for every
$0\le\varepsilon<\delta$,
$\max_j(b_j+\varepsilon s_j)=\max_j b_j+\varepsilon D$.
The actual parameter scales are positive and tend to zero because their
actual denominators are eventually $eN+c$ with $e>0$. Hence the displayed
identity applies eventually to THESE actual scales. The original maximum
identity at $q_N$ and $q_{\mathrm{vert}}$ gives the asserted exact equality for the
ORIGINAL $u$. The finite maximum is real, so extended-real conversion loses
nothing. Divide by the strictly positive actual scale to obtain the exact
eventual quotient; eventual constancy proves the limit.
\end{proof}

\paragraph{Subsequent weighted-degree boundary, not an input to this producer.}
For application to the original line, the global surface/SNC construction must
produce the actual original surface and selected actual central component
$C$, the actual stalk/parameter data above at every closed point of the
proper normal curve, and the actual original generic-field triangles.
It must identify all actual $q_{\mathrm{vert}}$ with one common original vertical point.
The canonical leading-residue field at each stalk must identify with the
SAME $K(C)$ and its actual local curve order, with its actual original $k$-residue-degree
weight on the central curve. These are geometric identifications, not desired slope or balance
certificates. The numerical part separately produces the ORIGINAL model
line's rational section and actual whole-point norm dictionary, regauged
with vertical order zero, and identifies its full degree with the genuine
weighted residual orders. The fixed-point producer assumes none of those conclusions.

Once those actual diagrams are derived, the weighted argument uses two
literal principal balances. For a common actual generic uniformizer, the
corrected order is the curve order of the actual leading $f_i$ minus
$(A_i/e)$ times the curve order of the actual leading $t_p$. Each has finite
support and whole residue-weighted sum zero by the genuine original
principal-degree theorem. Thus each $s_j$ has weighted sum zero. The SAME
active set is selected at the common original $q_{\mathrm{vert}}$; its maximum $D$ is
pointwise at least any one fixed active $s_j$. Nonnegative true residue
weights give $\sum_x[\kappa(x):k]D(x)\ge0$. The actual line-degree identity of the
numerical part then gives the desired nonnegative original
curve degree. Neither global equality nor positivity is silently supplied
to the active-profile theorem.
